% 工尺谱印本复刻 — 鹊桥仙·圣世呈符
% 参考: 示例/工尺谱/ref-2-upright.png
%
% 复刻范围 (基于 ref 可读部分):
%   鹊桥仙 一名广寒秋  (调名 + 别名)
%   圣世呈符 句  星垣辉朗鸿 韵
%   披香侍宴  上韵
%   金莲花炬照回廓  韵
%   上林游赏  韵    正院 ...
%   金鸡叫
%
% 工尺音字基本参考 ref 上方注释; 部分模糊处取近似. 板眼参考印本简化.
%
% 已知不一致:
%   - 句/韵 句读标记: 用 \小注 (本地宏 = {\small ...}) 近似, 印本中应是
%     专用的小字符号, 暂无独立命令
%   - 调名+别名 (鹊桥仙 一名广寒秋): 用普通文本 + {\small 一名广寒秋} 近似
%
\documentclass{ltc-guji}

\setmainfont{TW-Kai}

% 缩短列容量以匹配 ref-2 印本密度 (默认 21 字/列对此样过松)
\contentSetup{
    n-char-per-col = 12,
    n-column = 8,
}

\title{鵲橋仙}

\chapter{聖世呈符}

% 局部宏: 句读小标记 (一字大小约半个正文字)
\newcommand{\句读}[1]{{\small #1}}
\newcommand{\调名注}[1]{{\small #1}}

\begin{document}
\begin{正文}

% 调名 + 别名 (subtitle annotation)
鵲橋仙\空格\调名注{一名廣寒秋}

% --- 第一句: 聖世呈符·句 / 星垣輝朗鴻·韻 ---
\工尺{聖}{\音[板]{六}\音[眼]{凡}\音[板]{尺}}%
\工尺{世}{\音[眼]{凡}\音[板]{六}}%
\工尺{呈}{\音[眼]{尺}}%
\工尺{符}{\音[板]{凡}}%
\句读{句}

\工尺{星}{\音[板]{六}}%
\工尺{垣}{\音[眼]{凡}}%
\工尺{輝}{\音[板]{六}}%
\工尺{朗}{\音[眼]{尺}\音[板]{上}}%
\工尺{鴻}{\音[板]{凡}}%
\句读{韻}

% --- 第二句: 上 (单字) · 韻 ---
\工尺{上}{\音[板]{尺}}%
\句读{韻}

% --- 第三句: 披香侍宴 ---
\工尺{披}{\音[板]{工}}%
\工尺{香}{\音[眼]{工}}%
\工尺{侍}{\音[板]{六}}%
\工尺{宴}{\音[眼]{尺}}%

% --- 第四句: 金蓮花炬照回廓·韻 ---
\工尺{金}{\音[板]{五}\音[眼]{六}}%
\工尺{蓮}{\音[板]{五}}%
\工尺{花}{\音[眼]{六}}%
\工尺{炬}{\音[板]{五}}%
\工尺{照}{\音[眼]{尺}\音[板]{上}}%
\工尺{回}{\音[眼]{尺}}%
\工尺{廓}{\音[板]{凡}}%
\句读{韻}

% --- 第五句: 上林遊賞·韻 ---
\工尺{上}{\音[板]{六}}%
\工尺{林}{\音[眼]{凡}}%
\工尺{遊}{\音[板]{尺}\音[眼]{上}}%
\工尺{賞}{\音[眼]{凡}}%
\句读{韻}

% --- 第六句: 正院 (片段) ---
\工尺{正}{\音[板]{六}}%
\工尺{院}{\音[眼]{六}}%

% --- 第七句: 金雞叫 ---
\工尺{金}{\音[板]{凡}}%
\工尺{雞}{\音[眼]{六}}%
\工尺{叫}{\音[板]{尺}}%

\end{正文}
\end{document}
